% Front matter of the note's measurement rewrite (COMBINED-PLAN sec. 3(d)): the reader's guide, the what-is-measured box
% (it follows when cov2d's protocol rows arrive) and the definitions table. Owner: journal-editor. Nothing here is a new
% result: each entry restates a definition the note gives elsewhere and points to it. The reader's guide absorbs the
% executive summary's former Table 2 (tab:paper-note-map; move JE-03, history-cuts/move-manifest-je-20261117.tsv), with
% framing-map row A17's [W] wording: its caption, its first header and one row no longer name the methods paper.
\section*{Reader's guide}
\addcontentsline{toc}{section}{Reader's guide}
\label{sec:front-guide}

The limitations are collected in \S\ref{sec:limitations}, with the result component each one qualifies (Table~\ref{tab:result-status}). A reader checking the results can start there, then read the sixteen cells (\S\ref{sec:fived16}) and the two-dimensional result (\S\ref{sec:results}). The release package and the code are described in Apps.~\ref{app:release} and~\ref{app:codebase}. Choices made after results had been seen are listed in App.~\ref{app:devhistory}. How the note's numbers and statements are checked against the records is described in App.~\ref{app:verification}. Table~\ref{tab:paper-note-map} maps each topic to its sections.

Specific questions are answered in these tables: the sixteen cells and their per-cell uncertainties, Table~\ref{tab:fived16-cells}; the uncertainty groups and the correlations of their covariance, Tables~\ref{tab:fived16-groups} and~\ref{tab:fived16-corr}; the fake-data tests, Table~\ref{tab:fived16-fd}; the comparisons with generator predictions, Table~\ref{tab:fived16-gen}; and the decisions of the generator tests, Table~\ref{tab:joint5d}.

\begin{table}[H]
\centering\small
\caption{Where each topic is treated in this note.}
\label{tab:paper-note-map}
\begin{tabular}{@{}p{0.32\textwidth}p{0.64\textwidth}@{}}
\toprule
Topic & Note sections \\
\midrule
Introduction & \S\ref{sec:intro} \\
Data and observables & \S\ref{sec:experiment}, \S\ref{sec:method-binning} \\
Unfolding method and test design & \S\ref{sec:method}, \S\ref{sec:systematics}, \S\ref{sec:fived16-cov}; App.~\ref{app:statmethods}, App.~\ref{app:statmethods-nd}, App.~\ref{app:negweight} \\
Comparison with the published two-dimensional measurement & \S\ref{sec:results-compare} \\
Three- and four-observable central values & \S\ref{sec:3d} \\
Hadronic mass $W$ and the $(\Eavail,W)$ plane & \S\ref{sec:eavailw} \\
Recovery of controlled truth changes & \S\ref{sec:closure-map}, \S\ref{sec:fived16-fd}, \S\ref{sec:fived16-ibu}, \S\ref{sec:val-departure} \\
Reproducibility, stability and uncertainty calibration & \S\ref{sec:coverage}, \S\ref{sec:coverage-tunes}, \S\ref{sec:fived16-cov}, \S\ref{sec:fived16-fd}, \S\ref{sec:val-departure}, \S\ref{sec:gbdt-model-dependence}, \S\ref{sec:coverage-5d} \\
Generator tests & \S\ref{sec:joint5d} \\
Limitations & \S\ref{sec:limitations}; Tables~\ref{tab:test-index} and~\ref{tab:result-status} \\
Discussion and conclusions & \S\ref{sec:summary} \\
\bottomrule
\end{tabular}
\end{table}

